% !TEX program = xelatex
\documentclass[12pt,a4paper]{ctexart}

\usepackage{fontspec}
\setmainfont{Times New Roman}
\setCJKmainfont{SimSun}[BoldFont=SimHei,ItalicFont=KaiTi]
\setCJKsansfont{SimHei}
\setCJKmonofont{FangSong}

\usepackage[top=28mm,bottom=28mm,left=28mm,right=28mm]{geometry}
\usepackage{fancyhdr}
\usepackage{perpage}
\MakePerPage{footnote}

\setlength{\parindent}{2em}
\setlength{\parskip}{0pt}

\newcommand{\circledfootnotenumber}{%
  {\fontspec{SimSun}\ifcase\value{footnote}\or ①\or ②\or ③\or ④\or ⑤\or ⑥\or ⑦\or ⑧\or ⑨\else\arabic{footnote}\fi}}
\renewcommand{\thefootnote}{\circledfootnotenumber}

\pagestyle{fancy}
\fancyhf{}
\fancyfoot[C]{\thepage}
\renewcommand{\headrulewidth}{0pt}

\begin{document}

\begin{center}
    {\zihao{-3}\heiti 半个世纪后，我想对您说：我理解了道路\par}
    \vspace{0.8em}
    {\zihao{-4}\kaishu 姓名：\underline{\hspace{3.5cm}}\qquad 学号：\underline{\hspace{4cm}}\par}
\end{center}

\vspace{1em}

{\zihao{-4}\songti\setlength{\baselineskip}{20pt}\selectfont

毛泽东同志：

写这篇课程总结的时候，我先把题目看了很久。“半个世纪后，我想对您说”，这个说法比普通的论文题目更难下笔。它不是让我把教材内容重新概括一遍，也不是让我简单地表达敬意，而是要我把这一学期真正想清楚的一点东西说出来。2026年是您逝世五十周年。五十年听起来已经很远，但这门课让我感觉，许多历史问题其实并没有完全离开今天，只是换了场景、换了语言，继续要求后来的人去回答。

老实说，刚开始学习“毛泽东思想和中国特色社会主义理论体系概论”时，我对这门课的期待并不高。我以为它主要是概念、时间线和结论，认真听课、完成作业就可以了。后来我慢慢发现，如果只是这样学，确实很容易觉得它离自己很远。真正让我转变的是读《新民主主义论》。我第一次读的时候，最先注意到的是“中国向何处去”这句话。它不像一个已经知道答案的人写下的开场白，更像是一个时代被逼到路口时的追问。那时的中国到底应该往哪里走，不是靠热情就能解决，也不是把外国经验搬来就能解决。

读到这里，我才对“实事求是”有了比较具体的感觉。以前我觉得这个词太熟悉，熟悉到几乎不用再想。可是放在《新民主主义论》的语境里，它其实很艰难。您不是先给中国安排一个现成答案，而是先分析近代中国的社会性质、革命任务和各种力量，再说明为什么中国革命必须走新民主主义道路。\footnote{毛泽东：《新民主主义论》，《毛泽东选集》第2卷，北京：人民出版社，1991年，第662--706页。}这让我意识到，理论最有力量的时候，往往不是它说得多完整，而是它能够进入现实，把混乱的现实讲清楚。对我来说，这是这门课最重要的收获之一。

这学期课堂上还有一个内容让我印象很深，就是关于您的历史功绩和晚年探索曲折的讨论。以前我理解历史人物，常常会不自觉地把他们放进比较简单的评价框架里，要么只看功绩，要么只看问题。但课堂提醒我，历史不是这么平的。没有您和那一代共产党人的探索，中国不可能完成民族独立和人民解放，也不可能在一个贫弱破碎的基础上建立新的国家制度。新中国成立以后，工业体系、国防体系和社会主义基本制度的建立，也为后来国家发展留下了基础。可是，社会主义建设中的失误同样不能回避。把这些放在一起看，会比单纯赞颂更沉重，但也更真实。

我想，这种真实感正是学习历史的意义。我们纪念您，不应该只是把历史变成几句熟悉的话，而应该理解那些话背后的处境。比如“独立自主”，如果只从字面看，好像是一个很坚定的姿态；但放回当时的中国，它其实意味着一个落后大国必须在外部压力和内部困难中自己找路。比如“群众路线”，如果只背“从群众中来，到群众中去”，很容易背完就过去了；但读《湖南农民运动考察报告》时，我会感觉到，您对农民运动的判断来自对社会现场的观察，而不是从书本里直接推出的。\footnote{毛泽东：《湖南农民运动考察报告》，《毛泽东选集》第1卷，北京：人民出版社，1991年，第12--44页。}这让我对“人民”这个词多了一点敬畏。人民不是文章里用来增强气势的词，而是历史真正发生的地方。

这也让我想到今天的青年。我们当然不处在战争年代，也不需要用当年的方式投身革命。但如果说每一代青年都有自己的时代问题，那么今天的问题同样不轻。科技竞争、就业压力、产业升级、国家统一、文化传承，这些词平时看起来很大，好像离普通学生有距离，可是它们最后又会落回每个人的学习和选择里。以前我很容易把专业课、考试、作业和国家发展分开看，觉得前者是自己的事，后者是新闻里的事。学完这门课后，我不敢说自己已经有多高的认识，但至少会多想一步：我现在怎样学习，未来怎样工作，其实也会在某种程度上参与这个时代。

读《青年运动的方向》时，我也有类似感受。文章里谈青年要同工农群众结合，放到今天看，具体形式已经不同了，但它仍然提醒我，青年不能只在自己的小圈子里理解世界。\footnote{毛泽东：《青年运动的方向》，《毛泽东选集》第2卷，北京：人民出版社，1991年，第559--565页。}大学生活有时会让人把注意力集中在绩点、证书、升学和就业上，这些当然重要，可如果完全只围绕个人竞争打转，就很容易对社会现实变得迟钝。课程中的阅读和讨论至少让我意识到，所谓青年责任不是一句突然变大的口号，而是在日常学习里慢慢形成一种关心现实的习惯。

我也想诚实地说，这门课并没有让我一下子变得特别坚定或特别成熟。很多内容我仍然只是刚刚理解一点，有些历史问题也还需要继续读书才说得清楚。但它改变了我看问题的习惯。过去遇到复杂问题，我常常希望马上找到一个简单判断；现在我会觉得，先弄清楚问题从哪里来、和哪些条件有关，可能比急着表态更重要。您当年反复强调从中国实际出发，对今天的学生来说，仍然是一种很有用的提醒：不要只套用别人的答案，也不要把现实看得过于简单。

半个世纪后，我想对您说，我从这门课里重新理解了“道路”二字。道路不是写在纸上的路线图，也不是一开始就平坦清楚的选择。它是在矛盾、困难、尝试和修正中走出来的。中国革命如此，中国社会主义建设如此，今天中国式现代化的推进也是如此。对我个人来说，理解这一点以后，再看自己的学习和未来，就不会只想着怎样完成眼前任务，而会更愿意把自己放进一个更长的时间里去考虑。

如果这是一封迟到五十年的信，我最后想说的是：作为后来的一名普通学生，我未必能完全理解您所处时代的压力，也还没有能力回答今天中国面对的所有问题。但经过这一学期，我至少明白了一点，真正的纪念不是把历史人物放得很远，而是在自己的时代里继续学习怎样认识中国、怎样理解人民、怎样把个人的路同国家的路联系起来。这个答案还很朴素，却是我现在最真实的想法。

\vspace{1em}
\noindent{\heiti 参考文献}

\noindent [1] 毛泽东：《新民主主义论》，《毛泽东选集》第2卷，北京：人民出版社，1991年。

\noindent [2] 毛泽东：《湖南农民运动考察报告》，《毛泽东选集》第1卷，北京：人民出版社，1991年。

\noindent [3] 毛泽东：《青年运动的方向》，《毛泽东选集》第2卷，北京：人民出版社，1991年。

}

\end{document}
